\documentclass[10pt,a4paper]{amsart}
\usepackage[margin=19mm]{geometry}
\usepackage{amsmath,amssymb,mathtools}
\usepackage[scr=rsfs,cal=euler]{mathalfa}
\usepackage{xfrac}
\usepackage[unicode,hidelinks,bookmarks=false]{hyperref}
\newcommand{\ie}{i.e.\ }
\newcommand{\cf}{cf.\ }
\newcommand{\resp}{respectively\ }
\newcommand{\Subsection}{\subsection}
\providecommand{\nobreakdash}{\nobreak\hbox{-}}
\setlength{\emergencystretch}{2em}
\multlinegap0pt
\usepackage{bm}
\usepackage{bbm}
\usepackage{dsfont}
\usepackage{xspace}
\usepackage{cancel}
\usepackage{mathtools}
\usepackage{amsthm}
\makeatletter
\@ifundefined{definition}{\newtheorem{definition}{Definition}}{}
\@ifundefined{prop}{\newtheorem{prop}[definition]{Proposition}}{}
\makeatother
\makeatletter
\expandafter\ifx\csname dmc\endcsname\relax
\newenvironment{dmc}[1]{\renewcommand{\Title}{#1}\begin{defititle}} {\end{defititle}}
\fi
\newcommand{\e}{e}
\makeatother
\title[Obstructions to lifting quaternionic torus actions]{Obstructions to lifting quaternionic torus actions}
\author{Panagiotis Batakidis, Ioannis Gkeneralis}
\date{Source: arXiv:2608.07597v1 (CC BY 4.0)}
\begin{document}
\maketitle

\noindent\emph{Source and licence.} Obstructions to lifting quaternionic torus actions. Version arXiv:2608.07597v1 (CC BY 4.0), distributed under the Creative Commons Attribution 4.0 International licence, as recorded in the arXiv licence metadata for this exact version and shown on the abstract page https://arxiv.org/abs/2608.07597v1. Full rights notice: licenses/arxiv-2608.07597-CC-BY-4.0.txt.\par\smallskip

\noindent\textit{Editorial scope.} Counted unit: the Prop whose source label is \texttt{prop:local-lifting-chartwise} in the article (source0.tex lines 2503--2552, source label \texttt{prop:local-lifting-chartwise}), delivered together with the article's own complete original proof of it (source0.tex lines 2554--2877, one \texttt{proof} environment of 324 lines). No mathematics is written, repaired or re-derived: statement and proof are the article's own text line for line. Required context retained uncounted: def:local-Q-lifting (lines 2463--2482); eq:lifted-semidir-compatibility (lines 2474--2481). Every definition, display and label of those lines is retained unchanged. Omitted from this package and not redistributed: 1--2462, 2482--2502, 2553--2553, 2878--3601. Nothing inside a retained excerpt is omitted or altered: every retained body is delivered exactly as the article's lines are written, so no omission receipt of any kind accompanies this package. Citations: the retained excerpts carry no citation command at all, so no bibliography entry of the article is transcribed and no reference is redistributed. Reference adaptation: none was needed and none was made; every reference inside the delivered unit resolves inside it, so no unresolved reference or citation is produced. Printed slips kept literal: no printed word, symbol, hypothesis or number of the counted statement or of its proof was altered. Layout and class: the article's own class is \texttt{article} with packages including etex, inputenc, babel, soul, lmodern, textcomp, verbatim, appendix, cite, amsmath, amsthm, amsfonts, amssymb, mathtools; the wrapper is the lane's pinned \texttt{amsart} editorial wrapper at 10pt on A4, which restates the theorem-like environments the retained text needs and adds the article's own macro definitions that the retained text needs (\texttt{\textbackslash e}), with extra wrapper packages (bm, bbm, dsfont, xspace, cancel, mathtools). The delivered numbering is the unit's own. Assets: no retained span contains a figure environment, a graphics-inclusion command, a PDF-page inclusion, a table or a TikZ picture; no image file is copied into this package, the package declares no asset dependency and no image file is redistributed with it. Licence: version arXiv:2608.07597v1 of this article is distributed under the Creative Commons Attribution 4.0 International licence, as recorded in the arXiv licence metadata for this exact version and shown on the abstract page https://arxiv.org/abs/2608.07597v1. A full rights notice is delivered at licenses/arxiv-2608.07597-CC-BY-4.0.txt. Characters: every retained excerpt is pure ASCII, so the delivered engine, XeLaTeX with its default Latin Modern text font, reports no missing character.
\begin{definition}
\label{def:local-Q-lifting}
A \emph{lift to $P$ of the local \(Q^n\)-action \(\mathcal Q\) on M} is a lift of the globalized \(Q^n\)-action to \(\widetilde P\), namely a continuous
homomorphism
\[
\widetilde\phi_Q:
Q^n\longrightarrow
\operatorname{Isom}(\widetilde P),
\]
covering $\phi_Q:Q^n\longrightarrow\operatorname{Homeo}(\widetilde M)$,
such that it is compatible with the lifted deck transformations:
\begin{equation}
\label{eq:lifted-semidir-compatibility}
\widetilde\phi_\pi(a)\circ
\widetilde\phi_Q(u)\circ
\widetilde\phi_\pi(a)^{-1}
=
\widetilde\phi_Q\bigl(\rho(a)(u)\bigr),\;\;\forall a\in\pi_1(B_M),\;u\in Q^n.
\end{equation}
\end{definition}

\begin{prop}
\label{prop:local-lifting-chartwise}
The principal \(Q^k\)-bundle $Q^k\longrightarrow P\longrightarrow M$ admits a lift of the local \(Q^n\)-action \(\mathcal Q\) if and only if there exists a family
\[
\left\{
(P_\alpha,\widetilde\phi_\alpha,\varphi_\alpha^P)
\right\}_{\alpha\in A}
\]
with the following properties:
\begin{enumerate}
\item[(1)]
\[
Q^k\longrightarrow
P_\alpha
\overset{\pi_\alpha}{\longrightarrow}
\varphi_\alpha^M(U_\alpha^M)
\]
is a locally trivial principal \(Q^k\)-bundle. Thus, for every point \(y\in \varphi_\alpha^M(U_\alpha^M),\) there exists an open neighborhood \(V\subseteq \varphi_\alpha^M(U_\alpha^M)\) and a principal \(Q^k\)-bundle trivialization \(P_\alpha|_V\cong V\times Q^k\).

\item[(2)]
\[
\widetilde\phi_\alpha:
Q^n\longrightarrow
\operatorname{Isom}(P_\alpha)
\]
is a lift of the restriction of the regular \(Q^n\)-action to
\(\varphi_\alpha^M(U_\alpha^M)\);

\item[(3)]
\[
\varphi_\alpha^P:
P|_{U_\alpha^M}
\longrightarrow
P_\alpha
\]
is a principal \(Q^k\)-bundle isomorphism covering
\(\varphi_\alpha^M\);

\item[(4)] On every nonempty overlap \(U_{\alpha\beta}^M\), the induced bundle isomorphism \(\varphi_{\alpha\beta}^P:=\varphi_\alpha^P\circ(\varphi_\beta^P)^{-1}\) between the restricted principal \(Q^k\)-bundles over the corresponding overlap charts is \(\rho_{\alpha\beta}\)-equivariant with respect to the local lifts, that is,
\begin{equation}
\label{eq:local-lift-overlap}
\varphi_{\alpha\beta}^P\circ
\widetilde\phi_\beta(u)
=
\widetilde\phi_\alpha\bigl(\rho_{\alpha\beta}(u)\bigr)
\circ
\varphi_{\alpha\beta}^P,\;\forall u\in Q^n.
\end{equation}
\end{enumerate}
\end{prop}

\begin{proof}
Suppose first that a family satisfying the conditions above, exists. Let $u\in Q^n,\;
(\widetilde b,p_0)\in\widetilde P$
and assume that
\[
\widetilde b\in p^{-1}(U_\alpha^B),
\qquad
\widetilde\varphi_\alpha^B(\widetilde b)
=
\bigl(p(\widetilde b),a_\alpha\bigr).
\]
Define
\begin{equation}
\label{eq:local-lift-from-charts}
\widetilde\phi_Q(u)(\widetilde b,p_0)
:=
\left(
\widetilde b,\,
(\varphi_\alpha^P)^{-1}
\left(
\widetilde\phi_\alpha
\bigl(
(\rho_\alpha\circ\rho(a_\alpha))(u)
\bigr)
\bigl(\varphi_\alpha^P(p_0)\bigr)
\right)
\right).
\end{equation}
We first verify that the definition is independent of the choice of the local chart. Suppose that
\[
\widetilde b \in p^{-1}(U_\alpha^B)\cap p^{-1}(U_\beta^B),
\]
and write
\[
\widetilde\varphi_\alpha^B(\widetilde b) = \bigl(p(\widetilde b),a_\alpha\bigr), \qquad \widetilde\varphi_\beta^B(\widetilde b) = \bigl(p(\widetilde b),a_\beta\bigr).
\]
Then $a_\alpha=a_{\alpha\beta}a_\beta$.
Put $v := (\rho_\beta\circ\rho(a_\beta))(u)$.
By the overlap compatibility condition~\eqref{eq:local-lift-overlap},
\[
\varphi_{\alpha\beta}^P\circ \widetilde\phi_\beta(v) = \widetilde\phi_\alpha\bigl(\rho_{\alpha\beta}(v)\bigr) \circ \varphi_{\alpha\beta}^P.
\]
Moreover,
\begin{align*}
\rho_{\alpha\beta}(v)
&=
\bigl(
\rho_\alpha\circ
\rho(a_{\alpha\beta})\circ
\rho_\beta^{-1}
\bigr)
\bigl(
(\rho_\beta\circ\rho(a_\beta))(u)
\bigr)
\\
&=
\bigl(
\rho_\alpha\circ
\rho(a_{\alpha\beta}a_\beta)
\bigr)(u)
\\
&=
\bigl(
\rho_\alpha\circ\rho(a_\alpha)
\bigr)(u).
\end{align*}
Since $\varphi_{\alpha\beta}^P
=
\varphi_\alpha^P\circ(\varphi_\beta^P)^{-1}$, it follows that
\[(\varphi_\alpha^P)^{-1}
\left(
\widetilde\phi_\alpha
\bigl(
(\rho_\alpha\circ\rho(a_\alpha))(u)
\bigr)
\bigl(\varphi_\alpha^P(p_0)\bigr)
\right)=
(\varphi_\beta^P)^{-1}
\left(
\widetilde\phi_\beta
\bigl(
(\rho_\beta\circ\rho(a_\beta))(u)
\bigr)
\bigl(\varphi_\beta^P(p_0)\bigr)
\right).\]
Hence~\eqref{eq:local-lift-from-charts} is independent of the choice of \(\alpha\) and defines a lift of the global \(Q^n\)-action on \(\widetilde M\).

We next verify compatibility with the deck transformations. Let $a\in\pi_1(B_M),\;
u\in Q^n$,
and suppose that $\widetilde\varphi_\alpha^B(\widetilde b)
=
\bigl(p(\widetilde b),a_\alpha\bigr)$.
With our convention for the right deck action,
\[
\widetilde\varphi_\alpha^B(\widetilde b\cdot a)
=
\bigl(p(\widetilde b),a_\alpha a\bigr)
\]
and hence
\[
\widetilde\varphi_\alpha^B(\widetilde b\cdot a^{-1})
=
\bigl(p(\widetilde b),a_\alpha a^{-1}\bigr).
\]
Applying successively
\(\widetilde\phi_\pi(a)^{-1}\),
\(\widetilde\phi_Q(u)\), and
\(\widetilde\phi_\pi(a)\), we obtain
\begin{align*}
&
\widetilde\phi_\pi(a)\circ
\widetilde\phi_Q(u)\circ
\widetilde\phi_\pi(a)^{-1}
(\widetilde b,p_0)
\\
&=
\left(
\widetilde b,\,
(\varphi_\alpha^P)^{-1}
\left(
\widetilde\phi_\alpha
\bigl(
(\rho_\alpha\circ\rho(a_\alpha a))(u)
\bigr)
\bigl(\varphi_\alpha^P(p_0)\bigr)
\right)
\right)
\\
&=
\left(
\widetilde b,\,
(\varphi_\alpha^P)^{-1}
\left(
\widetilde\phi_\alpha
\bigl(
(\rho_\alpha\circ\rho(a_\alpha))
(\rho(a)(u))
\bigr)
\bigl(\varphi_\alpha^P(p_0)\bigr)
\right)
\right)
\\
&=
\widetilde\phi_Q\bigl(\rho(a)(u)\bigr)
(\widetilde b,p_0).
\end{align*}

Hence \(\widetilde\phi_Q\) is a lift of the local action in the sense of Definition~\ref{def:local-Q-lifting}.

For the opposite direction, suppose that a lift
\[
\widetilde\phi_Q:
Q^n\longrightarrow\operatorname{Isom}(\widetilde P)
\]
satisfying~\eqref{eq:lifted-semidir-compatibility} is given. Put
\[
P_\alpha:=P|_{U_\alpha^M},
\qquad
\pi_\alpha:=\varphi_\alpha^M\circ\pi_P,
\qquad
\varphi_\alpha^P:=\operatorname{id}_{P|_{U_\alpha^M}}.
\]
For \(p_0\in P_\alpha\), and \(\e\in \pi_1(B_M)\) the identity element, 
\[
\left(
(\widetilde\varphi_\alpha^B)^{-1}
\bigl(\pi_M(\pi_P(p_0)),e\bigr),
p_0
\right)
\]
is a point of \(\widetilde P\). We define
\[
\widetilde\phi_\alpha:
Q^n\longrightarrow\operatorname{Isom}(P_\alpha)
\]
by requiring
\begin{align*}
&
\widetilde\phi_Q\bigl(\rho_\alpha^{-1}(u)\bigr)
\left(
(\widetilde\varphi_\alpha^B)^{-1}
\bigl(\pi_M(\pi_P(p_0)),e\bigr),
p_0
\right)
\\
&\qquad =
\left(
(\widetilde\varphi_\alpha^B)^{-1}
\bigl(\pi_M(\pi_P(p_0)),e\bigr),
\widetilde\phi_\alpha(u)(p_0)
\right).
\end{align*}
This definition is well defined because the restriction of the untwisted action to the subgroup $Q^n\cong Q^n\times\{e\}$
fixes the \(\widetilde B_M\)-coordinate. Hence
\(\widetilde\phi_Q(\rho_\alpha^{-1}(u))\) maps a point of the form
\[
\left(
(\widetilde\varphi_\alpha^B)^{-1}
\bigl(\pi_M(\pi_P(p_0)),e\bigr),
p_0
\right)
\]
to another point with the same first coordinate, and therefore uniquely determines \(\widetilde\phi_\alpha(u)(p_0)\). Since \(\widetilde\phi_Q\) is a homomorphism and \(\rho_\alpha^{-1}\) is an automorphism of \(Q^n\), we have
\[
\widetilde\phi_\alpha(uv)
=
\widetilde\phi_\alpha(u)\circ
\widetilde\phi_\alpha(v)
\]
for every \(u,v\in Q^n\). Moreover, each \(\widetilde\phi_\alpha(u)\) is a principal \(Q^k\)-bundle automorphism covering the regular \(Q^n\)-action on
\(\varphi_\alpha^M(U_\alpha^M)\). Thus
\[
\widetilde\phi_\alpha:
Q^n\longrightarrow\operatorname{Isom}(P_\alpha)
\]
is a lift of the local regular action. It remains to verify the compatibility on overlaps. Let
\[
p_0\in P|_{U_{\alpha\beta}^M},
\qquad
b:=\pi_M(\pi_P(p_0))\in B_M.
\]
Using the transition function of the universal covering, we have
\[
(\widetilde\varphi_\beta^B)^{-1}(b,e)
=
(\widetilde\varphi_\alpha^B)^{-1}(b,a_{\alpha\beta}).
\]
Equivalently,
\[
(\widetilde\varphi_\beta^B)^{-1}(b,e)
=
\phi_\pi(a_{\alpha\beta}^{-1})
\left(
(\widetilde\varphi_\alpha^B)^{-1}(b,e)
\right).
\]
Therefore, by the lifted compatibility relation
\eqref{eq:lifted-semidir-compatibility},
\begin{align*}
\left(
(\widetilde\varphi_\beta^B)^{-1}(b,e),
\widetilde\phi_\beta(u)(p_0)
\right)
&=
\widetilde\phi_Q\bigl(\rho_\beta^{-1}(u)\bigr)
\left(
(\widetilde\varphi_\beta^B)^{-1}(b,e),
p_0
\right)
\\
&=
\widetilde\phi_Q\bigl(\rho_\beta^{-1}(u)\bigr)
\widetilde\phi_\pi(a_{\alpha\beta}^{-1})
\left(
(\widetilde\varphi_\alpha^B)^{-1}(b,e),
p_0
\right)
\\
&=
\widetilde\phi_\pi(a_{\alpha\beta}^{-1})
\widetilde\phi_Q
\bigl(
\rho(a_{\alpha\beta})\circ\rho_\beta^{-1}(u)
\bigr)
\left(
(\widetilde\varphi_\alpha^B)^{-1}(b,e),
p_0
\right).
\end{align*}
Since $\rho_{\alpha\beta}
=
\rho_\alpha\circ
\rho(a_{\alpha\beta})\circ
\rho_\beta^{-1}$,
we have $\rho(a_{\alpha\beta})\circ\rho_\beta^{-1}
=
\rho_\alpha^{-1}\circ\rho_{\alpha\beta}$ and hence the last expression can be computed as
\begin{align*}
\left(
(\widetilde\varphi_\beta^B)^{-1}(b,e),
\widetilde\phi_\beta(u)(p_0)
\right)
&=
\widetilde\phi_\pi(a_{\alpha\beta}^{-1})
\widetilde\phi_Q
\bigl(
\rho_\alpha^{-1}(\rho_{\alpha\beta}(u))
\bigr)
\left(
(\widetilde\varphi_\alpha^B)^{-1}(b,e),
p_0
\right)
\\
&=
\left(
(\widetilde\varphi_\beta^B)^{-1}(b,e),
\widetilde\phi_\alpha
\bigl(\rho_{\alpha\beta}(u)\bigr)(p_0)
\right).
\end{align*}
Thus
\[
\widetilde\phi_\beta(u)(p_0)
=
\widetilde\phi_\alpha
\bigl(\rho_{\alpha\beta}(u)\bigr)(p_0).
\]
More generally, with the bundle identifications \(\varphi_\alpha^P\) and \(\varphi_\beta^P\) restored, this becomes
\[
\varphi_{\alpha\beta}^P\circ
\widetilde\phi_\beta(u)
=
\widetilde\phi_\alpha
\bigl(\rho_{\alpha\beta}(u)\bigr)
\circ
\varphi_{\alpha\beta}^P.
\]
Therefore the family
\[
\left\{
(P_\alpha,\widetilde\phi_\alpha,\varphi_\alpha^P) \right\}_{\alpha\in A}
\]
has all the required properties.
\end{proof}

\end{document}
